\documentclass[11pt,a4paper]{article}
\usepackage[T1]{fontenc}
\usepackage[utf8]{inputenc}
\usepackage{amsmath,amssymb,amsthm}
\usepackage[margin=1in]{geometry}
\usepackage[colorlinks=true,linkcolor=blue,urlcolor=blue]{hyperref}
\theoremstyle{plain}
\newtheorem{theorem}{Theorem}
\newtheorem{lemma}[theorem]{Lemma}
\newtheorem{proposition}[theorem]{Proposition}
\newtheorem{corollary}[theorem]{Corollary}
\theoremstyle{definition}
\newtheorem{remark}[theorem]{Remark}
\newcommand{\Chat}{\widehat{C}}

\title{The extra $(x+1)$ factor:\\
a proof of Karttunen's second conjecture on the\\
binary columns of the Fibonacci numbers}
\author{Adrian Perez Fontelles\\ \small Independent researcher}
\date{23 August 2026}
\begin{document}
\maketitle

\begin{abstract}
OEIS A036284 reads one period of the $n$-th binary column of the Fibonacci numbers as a
polynomial over $\mathrm{GF}(2)$. A.~Karttunen conjectured in November 1998 that this
polynomial is divisible by $(x^{3}+1)^{2^{n-1}-1}$ (Conjecture 1) and that ``there is also
one extra $(x^{1}+1)$ factor present'' (Conjecture 2). We prove Conjecture 2, that is that
$(x^{3}+1)^{2^{n-1}-1}(x+1)$ divides the $n$-th term. Equivalently, the $(x+1)$-adic
valuation of the column polynomial is at least $2^{n-1}$, one more than the valuation of
$(x^{3}+1)^{2^{n-1}-1}$. The proof is a carry-periodicity identity read on the factor $x+1$
instead of on $x^{2}+x+1$; the point is simply that $x+1$ does not divide $x^{2}+x+1$, so
the cancellation that costs one power on the $x^{2}+x+1$ side costs nothing here. The note
is self-contained: the identity is stated and proved in Section~3.
\end{abstract}

\section{The conjecture}

Let $F(0)=0$, $F(1)=1$, $F(k)=F(k-1)+F(k-2)$. The bits in the $n$-th column from the right
of the binary expansions of $F(0),F(1),F(2),\dots$ repeat with period $3\cdot2^{n}$, and
OEIS A036284 (A.~Karttunen, 1 November 1998, \emph{Periodic vertical binary vectors of
Fibonacci numbers}) records one period read as a binary number:
\[
  a(n)=\sum_{k=0}^{3\cdot2^{n}-1}\Bigl(\bigl\lfloor F(k)/2^{n}\bigr\rfloor\bmod2\Bigr)2^{k},
  \qquad a(0),\dots,a(3)=6,\;24,\;1440,\;5728448.
\]
Identifying an integer with the $\mathrm{GF}(2)$ polynomial whose coefficient of $x^{k}$ is
bit $k$, the entry carries

\begin{quote}
\textbf{Conjecture:} For $n\ge1$, each term $a(n)$, when considered as a
$\mathrm{GF}(2)[X]$-polynomial, is divisible by the $\mathrm{GF}(2)[X]$-polynomial
$(x^{3}+1)^{\mathrm{A000225}(n-1)}$. \dots\ \textbf{Conjecture 2:} there is also one extra
$(x^{1}+1)$ factor present, see A136384.
\end{quote}

with $\mathrm{A000225}(n-1)=2^{n-1}-1$. Here we prove Conjecture 2, in the form
\[
  (x^{3}+1)^{2^{n-1}-1}(x+1)\ \Bigm|\ a(n)\qquad(n\ge1).
\]
As of the ``Last modified'' line on the live entry (23 August 2026) both statements are
still recorded as unproven conjectures and no proof appears in the entry's links or
programs.

\section{Notation}

Fix $n\ge1$ and set
\[
  L=3\cdot2^{n},\qquad q(x)=x^{2}+x+1,\qquad x^{3}+1=(x+1)q(x).
\]
Note that $q(1)=1$, so $x+1\nmid q$; this small fact is the whole difference between the
two conjectures. For nonzero $P\in\mathrm{GF}(2)[x]$ and irreducible $f$, let $v_{f}(P)$ be
the exponent of $f$ in $P$. By the Frobenius endomorphism,
\begin{equation}\label{eq:frob}
  x^{3\cdot2^{m}}+1=(x^{3}+1)^{2^{m}}\qquad(m\ge0).
\end{equation}

Let $r_{k}=F(k)\bmod2^{n+1}$, let $b(k)=\lfloor r_{k}/2^{n}\rfloor\in\{0,1\}$ be the $n$-th
bit and $s_{k}=r_{k}\bmod2^{n}$ the low part, so $r_{k}=b(k)2^{n}+s_{k}$. Let
\[
  \gamma(k)=\bigl\lfloor(s_{k-1}+s_{k-2})/2^{n}\bigr\rfloor\in\{0,1\}
\]
be the carry into position $n$ when $F(k-1)$ and $F(k-2)$ are added, and put
\[
  A(x)=\sum_{k=0}^{L-1}b(k)x^{k},\qquad \Chat(x)=\sum_{k=0}^{L/2-1}\gamma(k)x^{k}.
\]
$A$ is the polynomial attached to $a(n)$.

\section{The identity}

\begin{lemma}\label{lem:pisano}
For $m\ge1$ the sequence $(F(k)\bmod2^{m})_{k\ge0}$ is purely periodic with least period
$3\cdot2^{m-1}$.
\end{lemma}

\begin{proof}
Pure periodicity holds because $F(k-2)=F(k)-F(k-1)$ makes the recurrence invertible modulo
any $m$; the value $\pi(2^{m})=3\cdot2^{m-1}$ of the Pisano period is classical, see
Wall~\cite{wall}.
\end{proof}

Thus $(r_{k})_{k}$ has period $L$, so $b$, $s$, $\gamma$ are functions on
$\mathbb{Z}/L\mathbb{Z}$ and $r_{k}\equiv r_{k-1}+r_{k-2}\pmod{2^{n+1}}$ holds for every
$k\in\mathbb{Z}/L\mathbb{Z}$.

\begin{lemma}\label{lem:id}
In $\mathrm{GF}(2)[x]$ there is an $M$ with
\begin{equation}\label{eq:main}
  q(x)A(x)=\Chat(x)\,(x^{3}+1)^{2^{n-1}}+M(x)\,(x^{3}+1)^{2^{n}}.
\end{equation}
\end{lemma}

\begin{proof}
Splitting each summand into high bit and low part,
\[
  r_{k-1}+r_{k-2}=\bigl(b(k-1)+b(k-2)+\gamma(k)\bigr)2^{n}
  +\bigl((s_{k-1}+s_{k-2})\bmod2^{n}\bigr),
\]
which is legitimate since $s_{k-1}+s_{k-2}<2^{n+1}$. Reducing modulo $2^{n+1}$ and
comparing with $r_{k}=b(k)2^{n}+s_{k}$, where $s_{k}<2^{n}$, gives
\[
  b(k)\equiv b(k-1)+b(k-2)+\gamma(k)\pmod2\qquad(k\in\mathbb{Z}/L\mathbb{Z}).
\]
Multiplying by $x^{k}$ and summing over $\mathbb{Z}/L\mathbb{Z}$, where the shift
$k\mapsto k-1$ becomes multiplication by $x$ modulo $x^{L}+1$, yields
$q(x)A(x)\equiv\Gamma(x)\pmod{x^{L}+1}$ with $\Gamma(x)=\sum_{k=0}^{L-1}\gamma(k)x^{k}$.
Now $\gamma(k)$ is a function of $s_{k-1},s_{k-2}$ alone, and $s_{j}=F(j)\bmod2^{n}$ has
period $3\cdot2^{n-1}=L/2$ by Lemma~\ref{lem:pisano} applied with $m=n$. So the coefficient
list of $\Gamma$ is two copies of that of $\Chat$:
\[
  \Gamma(x)=\Chat(x)\bigl(1+x^{L/2}\bigr)=\Chat(x)(x^{3}+1)^{2^{n-1}}
\]
by \eqref{eq:frob}. Since $x^{L}+1=(x^{3}+1)^{2^{n}}$, also by \eqref{eq:frob}, this is
\eqref{eq:main}.
\end{proof}

\section{The theorem}

\begin{theorem}\label{thm:main}
For every $n\ge1$, $v_{x+1}(A)\ge2^{n-1}$. Consequently $(x^{3}+1)^{2^{n-1}-1}(x+1)$
divides $A$.
\end{theorem}

\begin{proof}
Both terms on the right of \eqref{eq:main} are divisible by $(x^{3}+1)^{2^{n-1}}$ and
$v_{x+1}\bigl((x^{3}+1)^{m}\bigr)=m$, so $v_{x+1}(qA)\ge2^{n-1}$. Since $q(1)=1$ we have
$v_{x+1}(q)=0$, hence $v_{x+1}(A)=v_{x+1}(qA)\ge2^{n-1}$.

The same equation read with $v_{q}$ gives $v_{q}(qA)\ge2^{n-1}$, and $v_{q}(q)=1$, so
$v_{q}(A)\ge2^{n-1}-1$ (this is Conjecture 1). Writing
$(x^{3}+1)^{2^{n-1}-1}(x+1)=(x+1)^{2^{n-1}}q^{\,2^{n-1}-1}$ and comparing exponents,
$v_{x+1}(A)\ge2^{n-1}$ and $v_{q}(A)\ge2^{n-1}-1$ are exactly what is needed.
\end{proof}

\begin{corollary}
For every $n\ge1$, $a(n)$ is divisible in $\mathrm{GF}(2)[X]$ by
$(x^{3}+1)^{2^{n-1}-1}(x+1)$. This is Karttunen's Conjecture 2.
\end{corollary}

\begin{remark}[honest accounting]
Once the identity of Lemma~\ref{lem:id} is in hand, Conjecture 1 and Conjecture 2 are the
two valuations one can read off it, and Conjecture 2 costs one extra line: it is what the
$(x+1)$-side gives when no power is lost to the cancellation. The entry states them as two
separate conjectures, and what is settled here is the second.
\end{remark}

\begin{remark}[sharpness]
Equation \eqref{eq:main} gives the exact value $v_{x+1}(A)=2^{n-1}+v_{x+1}(\Chat)$ whenever
$v_{x+1}(\Chat)<2^{n-1}$. Computation gives $v_{x+1}(\Chat)=0$ for $n=1$ and for
$4\le n\le12$, so the bound of Theorem~\ref{thm:main} is attained there; at $n=2$ and $n=3$
one has $v_{x+1}(\Chat)=1$ and 3, so $v_{x+1}(A)=3$ and 7 rather than 2 and 4. Whether
$v_{x+1}(\Chat)=0$ for all $n\ge4$ --- equivalently whether the number of carries in one
period of length $L/2$ is odd --- is not proved here.
\end{remark}

\section{Verification}

A verification script, written separately from this note, checks every assertion and prints
every number quoted. It (i) recomputes $a(0),\dots,a(3)$ and matches the published
$6,24,1440,5728448$; (ii) checks the carry identity of Lemma~\ref{lem:id} at all $L$
residues for $n=0,\dots,16$, that is up to $L=196608$, with no failure; (iii) checks the
half-period property of $\gamma$ over the same range, with no failure; (iv) verifies
\eqref{eq:main} by explicit polynomial division for $n=1,\dots,12$; (v) computes
$v_{x+1}(A)$, obtaining $2^{n-1}$ for $n=1$ and $4\le n\le12$ and the exceptional 3 and 7
at $n=2,3$, together with $v_{x+1}(\Chat)=0,1,3,0,\dots,0$; and (vi) divides $a(n)$ by
$(x^{3}+1)^{2^{n-1}-1}(x+1)$ directly for $n=1,\dots,12$, obtaining remainder zero in every
case.

\begin{thebibliography}{9}
\bibitem{wall} D.~D. Wall, \emph{Fibonacci series modulo $m$}, Amer. Math. Monthly 67
(1960), 525--532.
\bibitem{ln} R.~Lidl and H.~Niederreiter, \emph{Finite Fields}, 2nd ed., Cambridge
University Press, 1997.
\bibitem{oeis} N.~J.~A. Sloane et al., \emph{The On-Line Encyclopedia of Integer
Sequences}, \url{https://oeis.org/A036284}, entry A036284, consulted 23 August 2026.
\end{thebibliography}
\end{document}
